\chapter{Networking: GNAT.Sockets over the W5500}
\label{ch:networking}

Storage gave the application a filesystem; this chapter gives it a network. The
striking thing about the networking layer is the API the application sees: it is
the \emph{standard} \texttt{GNAT.Sockets} --- \texttt{Create\_Socket},
\texttt{Connect\_Socket}, \texttt{Send\_Socket}, a \texttt{Stream} over a socket.
Portable networking code compiles and runs here unchanged, within the subset this
chapter describes. Underneath that familiar surface sits a small stack: a
chip-neutral interface every NIC implements, and the WIZnet W5500 driver as the
first chip behind it.

\begin{figure}[ht]
\centering
\begin{tikzpicture}[node distance=2mm]
\node[layer, fill=blue!9] (app)  {Your application --- portable \texttt{GNAT.Sockets} code, \texttt{DNS\_Client}};
\node[layer, below=of app]  (fac)  {\texttt{GNAT.Sockets} facade --- interface registry, address types, the \texttt{Stream}};
\node[layer, below=of fac]  (intf) {\texttt{Net\_Devices.Device'Class} --- the chip-neutral contract (dispatched on)};
\node[layer, below=of intf] (wrap) {\texttt{ESP32S3.W5500.Net\_Device} --- the W5500 as the first backend};
\node[layer, below=of wrap] (eng)  {\texttt{ESP32S3.W5500.Sockets} --- the socket engine (\texttt{.Interrupts}, \texttt{.DHCP} alongside)};
\node[layer, below=of eng]  (xport){\texttt{ESP32S3.W5500} --- VDM frame transport + bring-up};
\node[layer, below=of xport] (spi) {\texttt{ESP32S3.SPI} --- task-safe SPI master $\rightarrow$ the chip};
\end{tikzpicture}
\caption{The networking stack. Everything above \texttt{Net\_Devices.Device'Class}
is chip-independent; everything below it is the W5500 driver. A second NIC slots in
as another backend without the layers above noticing.}
\label{fig:net-layers}
\end{figure}

\cnote{On an ESP-IDF system you reach for lwIP and the BSD socket API, with the
Wi-Fi or Ethernet \texttt{netif} underneath. The shape here is the same --- a
sockets API over a driver --- but the W5500 carries its \emph{own} TCP/IP stack in
silicon, so there is no lwIP to link: the ``stack'' is eight hardware sockets on
the chip, and our job is only to drive them and present the sockets API.}

\section{The W5500 and its driver}\index{W5500}\index{Ethernet}\index{SPI}
\label{ch:net:w5500}

The WIZnet W5500 is an Ethernet controller with a \emph{hardwired} TCP/IP stack:
ARP, ICMP, IP, TCP and UDP are implemented in the chip, which exposes eight
independent \emph{hardware sockets} and 32\,KB of buffer RAM, and is driven over
SPI. You do not run a protocol stack on the host; you write a destination address
and payload into a socket's registers and tell it to \texttt{SEND}. On this board
the chip is wired to SPI2 with \texttt{MISO=IO45}, \texttt{MOSI=IO4},
\texttt{SCLK=IO1}, \texttt{SCSn=IO39}, \texttt{INTn=IO3} (pull-up) and
\texttt{RSTn=IO11} (pull-up); the chip-select is an ordinary GPIO held low across
each frame, exactly as the SD-card driver does.

The driver is built in layers, each a child package, so the application can stop at
whichever level it needs:

\begin{itemize}[nosep]
\item \texttt{ESP32S3.W5500} --- the transport. The W5500's SPI ``variable data
length'' frame is a 16-bit address, a control byte (block-select + read/write), and
the data; this layer reads and writes registers and the socket buffers over the
task-safe SPI master, and brings the chip up (reset, \texttt{VERSIONR} probe, MAC /
IP / subnet / gateway, PHY link).
\item \texttt{ESP32S3.W5500.Sockets} --- the \emph{socket engine}. A self-contained
\texttt{Socket} handle with \texttt{Open\_TCP} / \texttt{Open\_UDP},
\texttt{Listen}, \texttt{Connect}, \texttt{Send} / \texttt{Receive},
\texttt{Send\_To} / \texttt{Receive\_From}, and a \texttt{Status} error model. This
is the level the higher layers build on.
\item \texttt{ESP32S3.W5500.Interrupts} --- optional. The single \texttt{INTn} line
(falling edge on \texttt{IO3}) wakes tasks blocked in \texttt{Wait\_Data} /
\texttt{Wait\_Connected} via per-socket suspension objects, so they sleep instead
of poll. Leaving it out costs nothing --- the engine polls.
\item \texttt{ESP32S3.W5500.DHCP} --- a DHCP client (see
\S\ref{ch:net:dns-dhcp}).
\item \texttt{ESP32S3.W5500.Net\_Device} --- the adapter that presents the chip as a
\texttt{GNAT.Sockets} interface (\S\ref{ch:net:netdevice}).
\end{itemize}

The engine and everything above it use controlled SPI handles and
\texttt{Ada.Real\_Time}, so the networking stack requires the \texttt{embedded} or
\texttt{full} profile.

\section{What GNAT.Sockets provides}\index{GNAT.Sockets}\index{socket}
\label{ch:net:gnatsockets}

\texttt{GNAT.Sockets} on a desktop is GNAT's binding to the operating system's
socket API. There is no operating system here, so this is a hand-written subset of
the \emph{same} package, backed by the W5500. Code written to it reads like any
sockets program:

\begin{lstlisting}
type Port_Type is range 0 .. 65535;
type Family_Type is (Family_Inet);                 --  IPv4
type Mode_Type   is (Socket_Stream, Socket_Datagram);

type Inet_Addr_Type is private;
function Inet_Addr (Image : String) return Inet_Addr_Type;   --  "a.b.c.d"
function Image (Value : Inet_Addr_Type) return String;
Any_Inet_Addr : constant Inet_Addr_Type;                     --  0.0.0.0

type Sock_Addr_Type is record
   Family : Family_Type := Family_Inet;
   Addr   : Inet_Addr_Type;
   Port   : Port_Type := 0;
end record;

type Socket_Type is private;
Socket_Error : exception;

procedure Create_Socket  (Socket : out Socket_Type; Family : Family_Type := Family_Inet;
                          Mode : Mode_Type := Socket_Stream);
procedure Bind_Socket    (Socket : in out Socket_Type; Address : Sock_Addr_Type);
procedure Listen_Socket  (Socket : in out Socket_Type; Length : Natural := 15);
procedure Accept_Socket  (Server : Socket_Type; Socket : out Socket_Type;
                          Address : out Sock_Addr_Type);
procedure Connect_Socket (Socket : in out Socket_Type; Server : Sock_Addr_Type);
procedure Send_Socket    (Socket : Socket_Type; Item : Stream_Element_Array;
                          Last : out Stream_Element_Offset; To : access Sock_Addr_Type := null);
procedure Receive_Socket (Socket : Socket_Type; Item : out Stream_Element_Array;
                          Last : out Stream_Element_Offset; From : access Sock_Addr_Type := null);
procedure Close_Socket   (Socket : in out Socket_Type);
function  Stream (Socket : Socket_Type) return Stream_Access;
\end{lstlisting}

A \texttt{Receive\_Timeout} socket option is provided through the standard
\texttt{Set\_Socket\_Option}, so a blocked receive can give up rather than hang:

\begin{lstlisting}
Set_Socket_Option (Sock, Socket_Level, (Receive_Timeout, Timeout => 3.0));
\end{lstlisting}

When the timeout elapses, \texttt{Receive\_Socket} raises \texttt{Socket\_Error},
just as it does on a desktop --- so portable code that already guards against that
works unchanged.

A few necessary departures from desktop \texttt{GNAT.Sockets}:

\begin{itemize}[nosep]
\item \textbf{IPv4 only} (\texttt{Family\_Inet}); the backend has no IPv6 yet.
\item \textbf{One bare-metal init.} Desktop \texttt{Initialize} takes no argument;
here you must say which chip backs the facade (\S\ref{ch:net:netdevice}).
\item \textbf{A finite number of sockets} --- the sum the interfaces provide (eight
per W5500). \texttt{Create\_Socket} past that raises \texttt{Socket\_Error}.
\item \textbf{Accept.} The W5500's listening socket \emph{becomes} the connection,
so \texttt{Accept\_Socket} returns the listener itself as the connected socket; to
accept another client, \texttt{Listen\_Socket} again.
\end{itemize}

\section{Wiring a driver in: the Net\_Device interface}\index{Net\_Devices}
\label{ch:net:netdevice}

The facade does not talk to the W5500 directly. Between them sits one small
package, \texttt{Net\_Devices}, defining the contract a network interface must
satisfy --- the part the facade actually needs. It is an Ada \texttt{interface}
(dispatched on at run time), plus chip-neutral address and status types:

\begin{lstlisting}
type Device is limited interface;
type Device_Access is access all Device'Class;

function Socket_Count (Self : Device) return Positive is abstract;
function Local_IP     (Self : Device) return IPv4_Address is abstract;
function Subnet_Mask  (Self : Device) return IPv4_Address is abstract;

procedure Open (Self : in out Device; Index : Natural; Mode : Transport;
                Local_Port : Port_Number; Result : out Status) is abstract;
procedure Close          (Self : in out Device; Index : Natural) is abstract;
procedure Listen         (Self : in out Device; Index : Natural; Result : out Status) is abstract;
procedure Wait_Connected (Self : in out Device; Index : Natural; Result : out Status) is abstract;
procedure Peer           (Self : in out Device; Index : Natural;
                          Addr : out IPv4_Address; Port : out Port_Number) is abstract;
procedure Connect        (Self : in out Device; Index : Natural; Host : IPv4_Address;
                          Port : Port_Number; Result : out Status) is abstract;
procedure Wait_Data      (Self : in out Device; Index : Natural; Result : out Status) is abstract;
procedure Send           (Self : in out Device; Index : Natural; Data : Stream_Element_Array;
                          Sent : out Natural; Result : out Status) is abstract;
procedure Receive        (Self : in out Device; Index : Natural; Into : out Stream_Element_Array;
                          Count : out Natural; Result : out Status) is abstract;
procedure Send_To        (Self : in out Device; Index : Natural; Host : IPv4_Address;
                          Port : Port_Number; Data : Stream_Element_Array; Result : out Status) is abstract;
procedure Receive_From   (Self : in out Device; Index : Natural; From : out IPv4_Address;
                          From_Port : out Port_Number; Into : out Stream_Element_Array;
                          Count : out Natural; Result : out Status) is abstract;
procedure Set_Receive_Timeout (Self : in out Device; Index : Natural; To : Duration) is abstract;
\end{lstlisting}

Sockets are addressed by an \texttt{Index} the device maps to its own per-socket
state (the W5500's eight hardware sockets). This is the \emph{offloaded-stack}
model: the device hands out TCP/UDP sockets directly. A raw-MAC chip with no stack
of its own would instead drive a software TCP/IP stack that implements this same
interface.

\subsection{The W5500 as a backend}

\texttt{ESP32S3.W5500.Net\_Device} is the first concrete \texttt{Net\_Devices.%
Device}. An \texttt{Instance} wraps one chip and owns that chip's eight
\texttt{Sockets.Socket} handles; each operation delegates to the engine and
converts the (identically-ordered) status and the address bytes. For the common
single-NIC board, one call wires everything together:

\begin{lstlisting}
type Instance is limited new Net_Devices.Device with private;
procedure Attach (Self : in out Instance; Dev : access ESP32S3.W5500.Device);
--  Single-NIC convenience: attach a built-in Instance and register it as default.
procedure Register_Default (Dev : access ESP32S3.W5500.Device);
\end{lstlisting}

The facade keeps a small \emph{registry} of interfaces and dispatches every socket
operation through it; a \texttt{Socket\_Type} records which interface it lives on
and which of that interface's sockets it is. Registration returns an id:

\begin{lstlisting}
type Interface_Id is range 0 .. 3;
function  Add_Interface (Device : Net_Devices.Device_Access) return Interface_Id;
procedure Initialize    (Device : Net_Devices.Device_Access);   --  register as default
\end{lstlisting}

So a board's bring-up ends with \texttt{ESP32S3.W5500.Net\_Device.Register\_Default
(Dev'Access)}, and from there the code is ordinary \texttt{GNAT.Sockets}. Because
the facade names only \texttt{Net\_Devices}, supporting a different chip means
writing one new backend (another implementation of the interface) and registering
it --- nothing above \texttt{Net\_Devices.Device'Class} changes, and a board may
register more than one. \emph{Per-destination routing across interfaces is not
wired yet: today every socket uses the default interface; the registry and the
per-socket interface id are the hooks a routing layer will use.}

\cnote{The interface plays the role of a Linux \texttt{net\_device} /
\texttt{struct net\_device\_ops}, or a \texttt{netif} driver in lwIP: a fixed set
of operations a chip implements so the protocol layer above stays generic. Here the
``protocol layer'' is the W5500's silicon; the interface is just how the facade
reaches it.}

\section{DNS the portable way, DHCP the chip way}\index{DNS}\index{DHCP}
\label{ch:net:dns-dhcp}

Two protocols sit just above sockets, and where they live tells you something about
the architecture.

\textbf{DNS} is an ordinary UDP request/response that happens \emph{after} the
interface is up, so it needs nothing chip-specific. It is written purely against
\texttt{GNAT.Sockets} and lives in a standalone, portable package --- the same
source would resolve names on a desktop:

\begin{lstlisting}
--  in DNS_Client
function Resolve (Server : Inet_Addr_Type; Name : String;
                  Addr : out Inet_Addr_Type;
                  Timeout : Duration := 0.0;
                  Local_Port : Port_Type := 13_001) return Boolean;
\end{lstlisting}

\textbf{DHCP} is the opposite kind of protocol: its job is to \emph{bootstrap} the
interface --- it must talk before an address exists and then program the obtained
IP, mask and gateway into the NIC. Both are operations outside the sockets API on
every platform (raw sockets and \texttt{ioctl} on a desktop; the W5500's
\texttt{Configure} registers here). So the DHCP client cannot be portable the way
DNS is; it rightly lives at the chip layer, \texttt{ESP32S3.W5500.DHCP}, built
straight on the socket engine. It runs the DORA exchange (Discover / Offer /
Request / Ack) and, on success, configures the chip. \texttt{Maintain} starts a
background task that acquires a lease and then keeps it --- renewing at $T1$,
rebinding at $T2$, re-acquiring on expiry:

\begin{lstlisting}
procedure Maintain (Dev : Sockets.Device_Access; MAC : MAC_Address;
                    Socket : Socket_Id := 0; On_Bound : Bound_Callback := null);
\end{lstlisting}

\section{The examples}
\label{ch:net:examples}

Six examples in \texttt{examples/} exercise the stack from the bottom up. Each
shares a tiny \texttt{W5500\_Dev} bring-up package (SPI + reset + a static address,
then \texttt{Register\_Default}); the DHCP example configures itself instead.

\subsection{\texttt{esp32s3\_w5500} --- TCP echo server}\index{TCP}

The full stack end to end: bring-up, the optional interrupts, and the facade. It
listens on \texttt{192.168.1.50:5000} and echoes each line back. This is the TCP
\emph{server} path --- \texttt{Listen\_Socket}, \texttt{Accept\_Socket} (which
reports the peer), \texttt{Receive\_Socket}, \texttt{Send\_Socket} --- with tasks
sleeping on \texttt{INTn} rather than polling.

\begin{lstlisting}
Create_Socket (Server, Family_Inet, Socket_Stream);
Bind_Socket   (Server, (Family_Inet, Any_Inet_Addr, 5000));
loop
   Listen_Socket (Server);
   Accept_Socket (Server, Conn, Peer);          --  blocks until a client connects
   loop
      Receive_Socket (Conn, Buf, Last);
      exit when Last < Buf'First;               --  peer closed
      Send_Socket (Conn, Buf (Buf'First .. Last), Sent);
   end loop;
   Close_Socket (Conn);
end loop;
\end{lstlisting}

\subsection{\texttt{esp32s3\_w5500\_http} --- HTTP GET (TCP client)}

The TCP \emph{client} path. It connects out to a web server, sends
\texttt{GET / HTTP/1.0}, and prints the response until the server closes ---
\texttt{Connect\_Socket} then a \texttt{Receive\_Socket} loop ending at
end-of-stream (\texttt{Last < Item'First}).

\subsection{\texttt{esp32s3\_w5500\_ntp} --- network time (UDP)}\index{NTP}\index{UDP}

A UDP datagram exchange: it sends a 48-byte NTP request to a time server with
\texttt{Send\_Socket (\dots, To => Server'Access)}, reads the reply with
\texttt{Receive\_Socket (\dots, From => \dots)}, takes the transmit timestamp, and
converts the NTP seconds to a civil date/time. Demonstrates \texttt{Socket\_%
Datagram} and the \texttt{To} / \texttt{From} addresses.

\subsection{\texttt{esp32s3\_w5500\_dns} --- name resolution}

A thin demonstration of \texttt{DNS\_Client.Resolve}: it resolves a hostname through
a resolver and prints the address. The protocol --- building the query, parsing the
answer with name-compression --- lives in the reusable module, so the example is a
few lines:

\begin{lstlisting}
if DNS_Client.Resolve (Inet_Addr ("8.8.8.8"), "example.com", Addr, Timeout => 3.0) then
   Put_Line ("example.com = " & Image (Addr));
end if;
\end{lstlisting}

\subsection{\texttt{esp32s3\_w5500\_dhcp} --- a self-configuring address}

No static IP: the board leases one. It calls \texttt{DHCP.Maintain} with a
library-level \texttt{On\_Bound} callback that prints each (re)bind, then the
background task acquires and keeps the lease for the life of the program. After the
first bind the chip is configured, so the socket engine and the facade are ready.

\section{Tying it together: the current weather}
\label{ch:net:weather}

\texttt{esp32s3\_w5500\_weather} uses the whole stack to answer a real question:
\emph{what is the weather at a latitude and longitude?} It fetches the current
conditions from Open-Meteo (\texttt{open-meteo.com}). The flow touches every layer
we have built --- UDP and TCP, name resolution, a receive timeout --- and is worth
reading start to finish.

Open-Meteo's forecast API answers over plain HTTP on port 80. That matters: the
W5500 has no TLS, so an HTTPS-only service would be unreachable without a software
crypto stack. Plain HTTP makes this a straight socket program.

\textbf{1. Bring up the interface.} \texttt{W5500\_Dev.Bring\_Up} does SPI + reset,
sets a static address, waits briefly for the PHY link, and registers the chip:
\texttt{ESP32S3.W5500.Net\_Device.Register\_Default (Dev'Access)}. From here it is
all \texttt{GNAT.Sockets}.

\textbf{2. Resolve the host (UDP).} The server's address is not hard-coded; it is
resolved at run time with the portable resolver, which sends a DNS query over a UDP
socket and parses the reply:

\begin{lstlisting}
if not DNS_Client.Resolve (DNS_Server, "api.open-meteo.com", Server_IP, Timeout => 5.0) then
   Put_Line ("[wx] DNS resolution failed");  return;
end if;
\end{lstlisting}

\textbf{3. Fetch the forecast (TCP).} Open a stream socket, connect to the resolved
address on port 80, and send a one-line HTTP request built from the latitude and
longitude; then read the whole response until the server closes:

\begin{lstlisting}
Create_Socket  (Sock, Family_Inet, Socket_Stream);
Connect_Socket (Sock, (Family_Inet, Server_IP, 80));
Send_Socket    (Sock, To_SEA (Request), Last);
loop
   Receive_Socket (Sock, Buf, Last);
   exit when Last < Buf'First;                  --  server closed the connection
   --  ... append Buf (Buf'First .. Last) to the response text ...
end loop;
Close_Socket (Sock);
\end{lstlisting}

where the request is

\begin{lstlisting}
"GET /v1/forecast?latitude=" & Latitude & "&longitude=" & Longitude
  & "&current_weather=true HTTP/1.0" & CRLF & "Host: api.open-meteo.com" & CRLF
  & "Connection: close" & CRLF & CRLF;
\end{lstlisting}

\textbf{4. Read the answer.} The reply is JSON. A small self-contained scrape (no
heap, no JSON library) pulls the fields out of the \texttt{current\_weather}
object --- anchoring past the \texttt{current\_weather\_units} object first, whose
same-named keys are unit \emph{strings} --- and a table turns the WMO weather code
into a word. The result, for Atlanta:

\begin{lstlisting}[style=sh]
[wx] resolving api.open-meteo.com ...
[wx] api.open-meteo.com = 188.40.99.226
[wx] forecast for 33.749, -84.388  (2026-06-25T20:45 UTC)
[wx]   temperature : 28.5 C
[wx]   wind        : 5.4 km/h from 184 deg
[wx]   conditions  : mainly clear
\end{lstlisting}

That single run exercises the entire chapter: a UDP exchange (DNS) and a TCP
client conversation (HTTP), each dispatched through \texttt{Net\_Devices.%
Device'Class} into the W5500 backend, down through the socket engine and the VDM
transport to the SPI master --- with a receive timeout guarding the wait, and the
W5500's silicon doing the TCP/IP. The application code, though, is just sockets:
the same program, pointed at a desktop \texttt{GNAT.Sockets}, would fetch the same
forecast.
